% GENERATED FILE. Edit canonical lessons, then run npm run generate:books.
% canonical-source-hash: fnv1a64:656d4ebe14e8c7bb
% canonical-lessons: JA-C300-kitsune, JA-C300-tanuki, JA-C300-bakeru, JA-C300-noroi, JA-C300-zurui

\chapter{The Fox That Changed Shape}
\label{ch:ja-the-fox-that-changed-shape}
\hlchaptermodality{300}

\begin{chapteropening}
\textbf{By the end of this chapter:} \emph{I can say what the fox and the badger did.}

Complete the last lesson of chapter 300: I can say what the fox and the badger did.
\end{chapteropening}

\section[\texorpdfstring{\ja{きつね} (kitsune)}{きつね (kitsune)}]{\ja{きつね} (kitsune) — a fox}

\label{lesson:JA-C300-kitsune}



\begin{quote}
Before the new one: say the Japanese for to protect, then the Japanese for a feudal lord.
\end{quote}

\subsection*{You'll want to know: \ja{きつね}}
\textbf{\ja{きつね}} — \emph{kitsune} — \textquotedblleft{}a fox\textquotedblright{}.

\begin{quote}
\textbf{\ja{やまに} \ja{きつねが} \ja{すんでいました。}} — \emph{yama ni kitsune ga sunde imashita} — \textquotedblleft{}A fox lived in the mountains.\textquotedblright{}
\end{quote}

\begin{cousinweb}
\textbf{\ja{きつね}} (\emph{kitsune}) is a native Japanese word (wago), inherited rather than borrowed from Chinese; its kanji was attached later as a kun'yomi (JMdict).
\end{cousinweb}

\begin{grammarlens}[title={What you've built}]
One more word for telling what happened.
\end{grammarlens}

\subsection*{Your turn}
\begin{itemize}
\raggedright
  \item \emph{Say it:} \emph{kitsune}
  \item \emph{Say it:} \emph{kitsune}, once more
  \item \emph{Say it:} \emph{yama ni kitsune ga sunde imashita} — A fox lived in the mountains.
  \item \emph{Recall:} say the Japanese for to protect, then the Japanese for a feudal lord, then say \emph{kitsune} again
\end{itemize}

\begin{tcolorbox}[breakable,colback=teal!4,colframe=teal!35!black,title={Before you move on}]
What does \ja{きつね} mean? (\textquotedblleft{}A fox\textquotedblright{}.) Say it once more.
\end{tcolorbox}
\section[\texorpdfstring{\ja{たぬき} (tanuki)}{たぬき (tanuki)}]{\ja{たぬき} (tanuki) — a raccoon dog}

\label{lesson:JA-C300-tanuki}



\begin{quote}
Before the new one: say the Japanese for a feudal lord, then the Japanese for a fox.
\end{quote}

\subsection*{You'll want to know: \ja{たぬき}}
\textbf{\ja{たぬき}} — \emph{tanuki} — \textquotedblleft{}a raccoon dog\textquotedblright{}.

\begin{quote}
\textbf{\ja{たぬきが} \ja{おじいさんの} \ja{うちに} \ja{きました。}} — \emph{tanuki ga ojiisan no uchi ni kimashita} — \textquotedblleft{}A raccoon dog came to the old man's house.\textquotedblright{}
\end{quote}

\begin{cousinweb}
\textbf{\ja{たぬき}} (\emph{tanuki}) is a native Japanese word (wago), inherited rather than borrowed from Chinese; its kanji was attached later as a kun'yomi (JMdict).
\end{cousinweb}

\begin{grammarlens}[title={What you've built}]
One more word for telling what happened.
\end{grammarlens}

\subsection*{Your turn}
\begin{itemize}
\raggedright
  \item \emph{Say it:} \emph{tanuki}
  \item \emph{Say it:} \emph{tanuki}, once more
  \item \emph{Say it:} \emph{tanuki ga ojiisan no uchi ni kimashita} — A raccoon dog came to the old man's house.
  \item \emph{Recall:} say the Japanese for a feudal lord, then the Japanese for a fox, then say \emph{tanuki} again
\end{itemize}

\begin{tcolorbox}[breakable,colback=teal!4,colframe=teal!35!black,title={Before you move on}]
What does \ja{たぬき} mean? (\textquotedblleft{}A raccoon dog\textquotedblright{}.) Say it once more.
\end{tcolorbox}
\section[\texorpdfstring{\ja{ばける} (bakeru)}{ばける (bakeru)}]{\ja{ばける} (bakeru) — to take the form of; to disguise oneself}

\label{lesson:JA-C300-bakeru}



\begin{quote}
Before the new one: say the Japanese for a fox, then the Japanese for a raccoon dog.
\end{quote}

\subsection*{You'll want to know: \ja{ばける}}
\textbf{\ja{ばける}} — \emph{bakeru} — \textquotedblleft{}to take the form of; to disguise oneself\textquotedblright{}.

\begin{quote}
\textbf{\ja{むかしばなしで、たぬきが} \ja{おじいさんに} \ja{ばけました。}} — \emph{mukashibanashi de, tanuki ga ojiisan ni bakemashita} — \textquotedblleft{}In an old tale, a raccoon dog turned itself into an old man.\textquotedblright{}
\end{quote}

Polite \textbf{\ja{ばけます}}, \emph{bakemasu}; an ichidan verb. Past (polite): \textbf{\ja{きつねが} \ja{おんなに} \ja{ばけました。}}, \emph{kitsune ga onna ni bakemashita}.

\begin{cousinweb}
\textbf{\ja{ばける}} (\emph{bakeru}) is a native Japanese word (wago), inherited rather than borrowed from Chinese; its kanji was attached later as a kun'yomi (JMdict).
\end{cousinweb}

\begin{grammarlens}[title={What you've built}]
One more word for telling what happened.
\end{grammarlens}

\subsection*{Your turn}
\begin{itemize}
\raggedright
  \item \emph{Say it:} \emph{bakeru}
  \item \emph{Say it:} \emph{bakeru}, once more
  \item \emph{Say it:} \emph{mukashibanashi de, tanuki ga ojiisan ni bakemashita} — In an old tale, a raccoon dog turned itself into an old man.
  \item \emph{Recall:} say the Japanese for a fox, then the Japanese for a raccoon dog, then say \emph{bakeru} again
\end{itemize}

\begin{tcolorbox}[breakable,colback=teal!4,colframe=teal!35!black,title={Before you move on}]
What does \ja{ばける} mean? (\textquotedblleft{}To take the form of; to disguise oneself\textquotedblright{}.) Say it once more.
\end{tcolorbox}
\section[\texorpdfstring{\ja{のろい} (noroi)}{のろい (noroi)}]{\ja{のろい} (noroi) — a curse; a spell}

\label{lesson:JA-C300-noroi}



\begin{quote}
Before the new one: say the Japanese for a raccoon dog, then the Japanese for to take the form of.
\end{quote}

\subsection*{You'll want to know: \ja{のろい}}
\textbf{\ja{のろい}} — \emph{noroi} — \textquotedblleft{}a curse; a spell\textquotedblright{}.

\begin{quote}
\textbf{\ja{おばあさんの} \ja{のろいで、おうさまは} \ja{ねこに} \ja{なりました。}} — \emph{obāsan no noroi de, ōsama wa neko ni narimashita} — \textquotedblleft{}The old woman's curse turned the king into a cat.\textquotedblright{}
\end{quote}

\begin{cousinweb}
\textbf{\ja{のろい}} (\emph{noroi}) is a native Japanese word (wago), inherited rather than borrowed from Chinese; its kanji was attached later as a kun'yomi (JMdict).
\end{cousinweb}

\begin{grammarlens}[title={What you've built}]
One more word for telling what happened.
\end{grammarlens}

\subsection*{Your turn}
\begin{itemize}
\raggedright
  \item \emph{Say it:} \emph{noroi}
  \item \emph{Say it:} \emph{noroi}, once more
  \item \emph{Say it:} \emph{obāsan no noroi de, ōsama wa neko ni narimashita} — The old woman's curse turned the king into a cat.
  \item \emph{Recall:} say the Japanese for a raccoon dog, then the Japanese for to take the form of, then say \emph{noroi} again
\end{itemize}

\begin{tcolorbox}[breakable,colback=teal!4,colframe=teal!35!black,title={Before you move on}]
What does \ja{のろい} mean? (\textquotedblleft{}A curse; a spell\textquotedblright{}.) Say it once more.
\end{tcolorbox}
\section[\texorpdfstring{\ja{ずるい} (zurui)}{ずるい (zurui)}]{\ja{ずるい} (zurui) — sly; cunning}

\label{lesson:JA-C300-zurui}



\begin{quote}
Before the new one: say the Japanese for to take the form of, then the Japanese for a curse.
\end{quote}

\subsection*{You'll want to know: \ja{ずるい}}
\textbf{\ja{ずるい}} — \emph{zurui} — \textquotedblleft{}sly; cunning\textquotedblright{}.

\begin{quote}
\textbf{\ja{ずるい} \ja{きつねは} \ja{たぬきを} \ja{だましました。}} — \emph{zurui kitsune wa tanuki o damashimashita} — \textquotedblleft{}The sly fox tricked the raccoon dog.\textquotedblright{}
\end{quote}

An i-adjective.

\begin{cousinweb}
\textbf{\ja{ずるい}} (\emph{zurui}) is a native Japanese word (wago), inherited rather than borrowed from Chinese; its kanji was attached later as a kun'yomi (JMdict).
\end{cousinweb}

\begin{grammarlens}[title={What you've built}]
One more word for telling what happened.
\end{grammarlens}

\subsection*{Your turn}
\begin{itemize}
\raggedright
  \item \emph{Say it:} \emph{zurui}
  \item \emph{Say it:} \emph{zurui}, once more
  \item \emph{Say it:} \emph{zurui kitsune wa tanuki o damashimashita} — The sly fox tricked the raccoon dog.
  \item \emph{Recall:} say the Japanese for to take the form of, then the Japanese for a curse, then say \emph{zurui} again
\end{itemize}

\begin{tcolorbox}[breakable,colback=teal!4,colframe=teal!35!black,title={Before you move on}]
What does \ja{ずるい} mean? (\textquotedblleft{}Sly; cunning\textquotedblright{}.) Say it once more.
\end{tcolorbox}
